
\section*{Notation}

We use $B(x,r)$ to denote the Euclidean (or $\ell_{2}$) ball centered
at $x$ with radius $r$: $\left\{ y:\norm{y-x}_{2}\le r\right\} $.
We use $\conv(X)$ to denote the convex hull of $X$, namely $\conv(X)=\{\sum\alpha_{i}x_{i}:\ \alpha_{i}\geq0,\sum\alpha_{i}=1,x_{i}\in X\}$.
We use $\left\langle x,y\right\rangle $ to denote the $\ell_{2}$
inner product $x^{T}y$ of $x$ and $y$. For any two points $x,y\in\Rn$,
we view them as column vectors, and use $[x,y]$ to denote $\conv(\{x,y\})$,
namely, the line segment between $x$ and $y$. 

We use $e_{i}$ to denote the coordinate vector whose $i$-th coordinate
is 1 and all other coordinates are $0$.


\section*{Functions}
\begin{defn}
For any $L\geq0$, a function $f:V\rightarrow W$ is $L$-Lipschitz
if $\norm{f(x)-f(y)}_{W}\leq L\norm{x-y}_{V}$ where the norms $\norm{\cdot}_{V}$
and $\norm{\cdot}_{W}$ are $\ell_{2}$ norms if unspecified.
\end{defn}

\begin{defn}
A function $f\in\mathcal{C}^{k}(\Rn)$ if $f$ is a $k$-differentiable
function and its $k^{\mathrm{th}}$ derivative is continuous.
\end{defn}

\begin{defn}
\label{def:lower-semi}A function ${\displaystyle f:X\subseteq\R^{n}\to\R}$
is called lower semi-continuous at a point $x_{0}\in X$ if for every
real $y<f(x_{0})$ there exists a neighborhood $U$ of $x_{0}$ such
that $f(x)>y$ for all $x\in U$. Equivalently, $f$ is lower semi-continuous
at $x_{0}$ iff 

\[
{\displaystyle \liminf_{x\to x_{0}}f(x)\geq f(x_{0})}.
\]
The function is lower semi-continuous if is it so at every point in
its domain. The definition of upper semi-continuous is analogous with
the inequalities reversed and $\liminf$ replaced by $\limsup$. 
\end{defn}

\begin{thm}[Taylor's Remainder Theorem]
\label{thm:taylor_expansion}For any $g\in\mathcal{C}^{k+1}(\R)$,
and any $x$ and $y$, there is a $\zeta\in[x,y]$ such that
\[
g(y)=\sum_{j=0}^{k}g^{(j)}(x)\frac{(y-x)^{j}}{j!}+g^{(k+1)}(\zeta)\frac{(y-x)^{k+1}}{(k+1)!}.
\]
\end{thm}

\begin{defn}
The $convolution$ of two functions $f,g:\R^{n}\rightarrow\R$ is
defined as $h(x)=\int_{\R^{n}}f(z)g(x-z)\,dz.$
\end{defn}


\section*{Linear Algebra}
\begin{defn}
\label{def:PSD}A real symmetric matrix $A$ is positive semi-definite
(PSD) if $x^{\top}Ax\geq0$ for all $x\in\Rn$. Equivalently, a real
symmetric matrix with all nonnegative eigenvalues is PSD. We write
$A\succeq0$ if $A$ is PSD and $A\succeq B$ if $A-B$ is PSD.
\end{defn}

\begin{defn}
For any matrix $A$, we define its trace, $\tr A=\sum A_{ii}$, Frobenius
norm, $\,\|A\|_{F}^{2}=\tr(A^{\top}A)=\sum_{i,j}A_{ij}^{2}$, and
operator norm, $\,\|A\|_{\op}=\sup_{\|x\|_{2}=1}\|Ax\|_{2}$. Note
that $x^{\top}Ax=\tr(Axx^{\top})$ and in general, $\tr(AB)=\tr(BA)$.
\end{defn}

For symmetric $A$, we have $\tr A=\sum_{i}\lambda_{i}$, $\|A\|_{F}^{2}=\sum_{i}\lambda_{i}^{2}$
and $\|A\|_{\op}=\max_{i}|\lambda_{i}|$ where $\lambda_{i}$ are
the eigenvalues of $A$.

For a vector $x$, $\norm x_{A}=\sqrt{x^{T}Ax}$; for a matrix $B$,
\[
\norm B_{A}=\sup_{x}\frac{\norm{Bx}_{A}}{\norm x_{A}}.
\]

\begin{lem}[Sherman-Morrison-Woodbury]
 For matrix $A\in\R^{n\times n}$ and vectors $u,v\in\R^{n}$, we
have
\[
\left(A+uv^{\top}\right)^{-1}=A^{-1}-\frac{A^{-1}uv^{\top}A^{-1}}{1+u^{\top}A^{-1}v}.
\]
For matrices $U,V\in\R^{n\times k},$we have
\[
\left(A+UV^{\top}\right)^{-1}=A^{-1}-A^{-1}U\left(I+V^{\top}A^{-1}U\right)^{-1}V^{\top}A^{-1}.
\]
\end{lem}


\section*{Probability}
\begin{defn}
The Total Variation (TV) or $\ell_{1}$-distance between two distributions
with measures $\rho,\nu$ with support $\Omega$ is 
\[
d_{TV}(\rho,\nu)=\frac{1}{2}\int_{\Omega}\left|\rho(x)-\nu(x)\right|\,dx=\sup_{S\subset\Omega}\left|\rho(S)-\nu(S)\right|=\sup_{S\subset\Omega}\rho(S)-\nu(S).
\]
\end{defn}

The following ``distances'' are not symmetric.
\begin{defn}
\label{defn:KL-divergence}The KL-divergence of a density $\rho$
with respect to another density $\nu$, also known as relative entropy
of $\rho$ with respect to $\nu$, is 
\[
\KL(\rho\|v)=\int\rho(x)\log\frac{\rho(x)}{\nu(x)}dx.
\]
\end{defn}

\begin{defn}
The $\chi$-squared distance of a densities $\rho$ with respect to
another density $\nu$ is defined as  
\[
\chi^{2}(\rho,\nu)=\E_{\nu}\left(\left(\frac{\rho(x)}{\nu(x)}-1\right)^{2}\right)=\E_{\rho}\left(\frac{\rho(x)}{\nu(x)}\right)-1.
\]
\end{defn}

\begin{defn}
\label{defn:The-Wasserstein--distance}The Wasserstein $p$-distance
between two probability measures $\rho,\nu$ over a metric space $M$
is defined as 
\[
W_{p}(\rho,\nu)=\left(\inf_{\gamma\in\Gamma(\rho,\nu)}\E_{x,y\sim\gamma}\left(d(x,y)^{p}\right)\right)^{1/p}
\]
where $\Gamma(\rho,\nu)$ is the set of all couplings of $\rho$ and
$\nu$ (joint probability measures with support $M\times M$ whose
marginals are $\rho$ and $\nu$).
\end{defn}

\begin{defn}
\label{def:p-dist} For $\mu,\nu$ probability measures in $\R^{n}$,
the \emph{$f$-divergence} of $\mu$ towards $\nu$ with $\mu\ll\nu$
is defined as, for a convex function $f:\R_{+}\to\R$ with $f(1)=0$
and $f'(\infty)=\infty$, 
\[
D_{f}(\mu\|\nu):=\int f\left(\frac{d\mu}{d\nu}\right)\,d\nu\,.
\]
For $q\in(1,\infty)$, the \emph{$\KL$-divergence and $\chi^{q}$-divergence}
correspond to $f(x)=x\log x$ and $x^{q}-1$, respectively. The \emph{$q$-R\'enyi
divergence} is defined as 
\[
\RR_{q}(\mu\|\nu):=\tfrac{1}{q-1}\log\left(\chi^{q}(\mu\|\nu)+1\right).
\]
The \emph{R\'enyi-infinity divergence} is defined as
\[
\RR_{\infty}(\mu\|\nu):=\log\esssup_{\mu}\frac{d\mu}{d\nu}\,.
\]
\end{defn}

\begin{fact}
$\KL=\lim_{q\downarrow1}\RR_{q}\leq\RR_{q}\leq\RR_{q'}\leq\RR_{\infty}$
for $1\leq q\leq q'$ and $2\,\norm{\cdot}_{\tv}^{2}\leq\KL\leq\RR_{2}=\log(\chi^{2}+1)\leq\chi^{2}$. 
\end{fact}

We refer readers to \cite{van2014renyi} for other basic properties
of  R\'enyi-divergence.

\begin{defn}
The \emph{marginal} of a distribution in $\R^{n}$ with density $\nu$
in the span of a $k$-dimensional subspace $V$ is defined as 
\[
g(x)=\int_{y\in V^{\perp}}\nu(x,y)\,dy
\]
for any $x\in V$.

Note that the marginal is a convolution of the density with the indicator
function of the subspace. 
\end{defn}

\begin{defn}
A distribution $D$ with support contained in $\R^{n}$ is said to
be \textbf{isotropic} if $\E_{D}(x)=0$ and $\E_{D}(xx^{\top})=I$. 
\end{defn}


\section*{Geometry}
\begin{defn}
We denote the unit Euclidean ball as $B^{n}=\left\{ x:\norm x_{2}\le1\right\} .$
More generally, we define the $\ell_{p}$-norm ball of radius $r$
centered at $z$ as $B_{p}^{n}(z,r)=\left\{ x:\norm{x-z}_{p}\le r\right\} $.
\end{defn}

\begin{defn}
The \emph{Minkowski sum }of two sets $A,B\subset\R^{n}$ is defined
as $A+B=\left\{ x+y:x\in A,y\in B\right\} $.
\end{defn}


